% !TEX root = ../../tree_notes.tex
% =====================================================================
%  5.4  Euler Tour for Distance and an RMQ Preview          sec:c24-dist
%  Exports: def:c24-tour  thm:c24-dist  prop:c24-rmqpreview
%           lst:c24-tour-build  ex:c24-dist-run
%           prob:c24-pathdir  lst:c24-pathdir
%           prob:c24-leafpairs lst:c24-leafpairs  practice: LC 1522 (N-ary diameter)
%  NOTE: LCA-by-Euler+RMQ is only PREVIEWED here; the proof and the
%  binary-lifting machinery live in the LCA chapter (\chref{c11}).
% =====================================================================
\section{Euler Tour for Distance and an RMQ Preview}\label{sec:c24-dist}

\begin{previouslybox}
\begin{enumerate}[nosep, leftmargin=1.6em]
  \item On the 2-tick clock, why does $u$ being an ancestor of $v$ read as
        $[\tin(v),\tout(v)) \subseteq [\tin(u),\tout(u))$? (\thmref{c24-anc})
  \item What is $\depth(v)$ for every node of $G$, and the root-path sums to nodes $4$ and
        $6$? (\secref{c24-path}, \figref{c24-running})
  \item Why did \secref{c24-flatten} hand off range \emph{answering} to the query chapters --
        and what does an RMQ have to do with LCA? (\chref{c13}, \chref{c11})
  \item State $\dist(u,v)$ in a tree as a count of edges on the unique $u$--$v$ path.
        (\chref{c01} unique-path property)
\end{enumerate}
\end{previouslybox}

\begin{topicintro}
The tour's final superpower: \emph{distance}. Record the DFS walk that revisits a node every
time it returns from a child -- a sequence $\mathcal{T}$ of length $2n-1$ -- and index it by
each node's \emph{first} appearance. The magic identity
$\dist(u,v) = \depth(u) + \depth(v) - 2\,\depth{\lca(u,v)}$ reduces distance to a lowest
common ancestor, and, astonishingly, $\lca(u,v)$ is just the \emph{shallowest} node in
$\mathcal{T}$ between the two first-appearances -- a \term{range minimum query} on the depth
array. That turns ``distance between two nodes'' into ``preprocess an array, answer RMQ in
$O(1)$''. We build the tour and the distance formula here; the RMQ-as-LCA \emph{equivalence},
its proof, and the faster binary-lifting alternative are the LCA chapter's to own
(\chref{c11}) -- this is the preview that motivates it.
\textbf{Prerequisites.} \secref{c24-times} ($\tin/\tout$, the 2-tick clock);
\secref{c24-path} (depth, root-paths); $\depth$ from \secref{c04-rooting}.

\textbf{After this section you can}
\begin{itemize}[nosep]
  \item build the $2n{-}1$ Euler tour and its first-occurrence array;      \hfill\lv{4}
  \item state and apply $\dist(u,v)=\depth u+\depth v-2\depth{\lca}$;          \hfill\lv{3}--\lv{4}
  \item see that $\lca$ $=$ an RMQ over the tour's depth array (the preview);   \hfill\lv{4}--\lv{5}
  \item solve root-path and leaf-pair distance problems by stitching downward paths. \hfill\lv{3}--\lv{4}
\end{itemize}
\end{topicintro}

\begin{recognitionbox}[Pattern card -- distance, LCA, and the tour]
\reckey{the question is about the \emph{path between two named nodes} -- their
\emph{distance}, the \emph{directions} to walk from one to the other, or a count over
\emph{leaf pairs within distance $d$} -- typically with many such pairs/queries.}
\recuse{
  \item ``distance between $u$ and $v$'' / ``is this pair of leaves close enough'' -- the
        $\dist$ formula via LCA (\thmref{c24-dist});
  \item ``print the moves $L/R/U$ from node $A$ to node $B$'' -- root-paths to their LCA
        (\probref{c24-pathdir});
  \item ``count pairs at a given distance'' -- combine leaf-depth lists at each common
        ancestor (\probref{c24-leafpairs});
  \item \emph{many} online distance/LCA queries -- reduce to RMQ on the tour
        (\propref{c24-rmqpreview}) and go build the structure in \chref{c11}.
}
\recnotuse{
  \item root-path-only aggregates -- \secref{c24-path}, no second endpoint;
  \item the LCA \emph{machinery} proved and coded -- that is \chref{c11}, not this preview;
  \item subtree ranges -- \secref{c24-flatten}.
}
One-line test: \emph{two arbitrary endpoints and a path between them? Distance is a LCA in
disguise -- tour $+$ $\dist$ formula, this section.}
\end{recognitionbox}

\subsection{Intuition}

Do the DFS, but \emph{log the current node every time you stand on it} -- on entry and again
each time you climb back from a child. For $G$ (order $0,1,4,5,2,3,6$) the walk of the
\emph{cursor} is
\[
  \mathcal{T} = [\,0,\,1,\,4,\,1,\,5,\,1,\,0,\,2,\,0,\,3,\,6,\,3,\,0\,],
\]
length $13 = 2\cdot 7 - 1$. Each node appears once per incident edge-plus-one, and the
\emph{first} time a node appears records when the DFS met it. Now look at the segment of
$\mathcal{T}$ between the first appearances of any two nodes $u, v$: the walk must climb from
$u$ up to their common ancestor and back down to $v$, and the \emph{shallowest} node it
touches in that window is exactly their lowest common ancestor.

\begin{figure}[H]
\centering
\begin{tikzpicture}[diagram, x=11mm, y=8mm]
  % tour as a polyline of depths
  \foreach \i/\dep in {0/0,1/1,2/2,3/1,4/2,5/1,6/0,7/1,8/0,9/1,10/2,11/1,12/0} {
     \node[circle, inner sep=1.6pt, fill=primaryblue!40] (p\i) at (\i,-\dep) {};
  }
  \draw[primaryblue, line join=round] (p0)--(p1)--(p2)--(p3)--(p4)--(p5)--(p6)--(p7)--(p8)--(p9)--(p10)--(p11)--(p12);
  % labels: node ids at each step
  \foreach \i/\lab in {0/0,1/1,2/4,3/1,4/5,5/1,6/0,7/2,8/0,9/3,10/6,11/3,12/0} {
     \node[annot, font=\tiny, below=0.5mm] at (\i,-0.05) {\lab};
  }
  % window first(4)=2 .. first(6)=10
  \draw[accentorange, line width=1.2pt, dashed] (1.6,0.6) -- (1.6,-2.6);
  \draw[accentorange, line width=1.2pt, dashed] (10.4,0.6) -- (10.4,-2.6);
  \node[annot, text=accentorange, font=\scriptsize, align=center] at (6,0.9)
    {window $[\mathrm{first}(4),\mathrm{first}(6)]$};
  \fill[warningred] (p6) circle (2.4pt);
  \node[annot, font=\scriptsize, text=warningred, align=left] at (6,-2.9)
    {shallowest node in the window $=0=\lca(4,6)$};
\end{tikzpicture}
\caption{The Euler tour of $G$ drawn as its depth profile (numbers under the dots are node
ids; the vertical position is $\depth$). Between the first appearances of $4$ and $6$
(dashed) the walk dips up to node $0$ -- the unique \emph{minimum depth} in that window --
which is $\lca(4,6)$. Distance is then $\depth(4)+\depth(6)-2\depth(0)$.}\label{fig:c24-tour}
\end{figure}

Read the depths of that window; its minimum (an RMQ) is the LCA. So \emph{one} $O(n)$ tour
plus \emph{any} RMQ structure (sparse table $\bigO{n\log n}$ preprocess, $O(1)$ query, or
the segment tree of \chref{c13}) answers distances in $O(1)$ -- the equivalence the LCA
chapter proves (\chref{c11}).

\subsection{Definitions and Notation}

\begin{definition}{Euler tour and first occurrence}{c24-tour}
The \term{Euler tour} (contour, $2n{-}1$ walk) of a rooted tree is the sequence
$\mathcal{T}$ obtained by recording the current node each time the DFS \emph{arrives} at it
and each time it \emph{returns} to it from a child. $|\mathcal{T}| = 2n-1$. The
\term{first-occurrence} array $\mathrm{first}[v]$ is the index of the earliest copy of $v$
in $\mathcal{T}$ (this is $\tin(v)$ under the entry-only preorder rank). The depth sequence
along the tour is $d[i] = \depth(\mathcal{T}[i])$.
\end{definition}

\noindent\textbf{In plain words.} Walk the outline of the tree with a pencil, jotting the
node every time the pencil is on it. Between two nodes' first jottings, the pencil necessarily
passes over their lowest shared point -- and that point is the shallowest jot in between.

\begin{theorem}{Distance via the LCA}{c24-dist}
For any nodes $u, v$ in a tree with a fixed root,
\[
   \dist(u, v) \;=\; \depth(u) + \depth(v) - 2\,\depth\!\big(\lca(u,v)\big).
\]
\end{theorem}
\begin{proofbox}
The unique $u$--$v$ path (Ch.~1) goes up from $u$ to $\ell = \lca(u,v)$ and down from $\ell$
to $v$, sharing only $\ell$. The $u$--$\ell$ segment has $\depth(u)-\depth(\ell)$ edges, the
$\ell$--$v$ segment $\depth(v)-\depth(\ell)$, and these are edge-disjoint, so
$\dist(u,v) = (\depth u - \depth \ell) + (\depth v - \depth \ell) = \depth u + \depth v -
2\depth\ell$.
\end{proofbox}

\begin{proposition}{LCA is a range minimum on the tour (preview)}{c24-rmqpreview}
$\lca(u,v)$ is the node of \emph{minimum} depth among
$\mathcal{T}[\min(\mathrm{first}\,u, \mathrm{first}\,v) \,..\, \max(\mathrm{first}\,u, \mathrm{first}\,v)]$.
Hence, with an RMQ structure over $d[\,]$, LCA -- and by \thmref{c24-dist} distance -- answers
in $O(1)$ after $\bigO{n}$ or $\bigO{n\log n}$ preprocessing.
\end{proposition}
\begin{proofbox}[Proof deferred to the LCA chapter]
Between first meeting $u$ and first meeting $v$, the DFS cursor must climb from $u$ up to
their common ancestors and back down to $v$; the set of nodes it visits in that window is
exactly the union of the two root-paths from $\ell$ to $u$ and $v$, so the shallowest is
$\ell$, and $\ell$ is strictly shallower than any node below it on either branch -- giving
uniqueness of the minimum \emph{node} (ties in depth never straddle two different nodes at the
same minimum here). A rigorous proof, the tie-handling, and the RMQ implementation (sparse
table / segment tree / $\pm$1-Cartesian $O(n)$) are the subject of \chref{c11}.
\end{proofbox}
\begin{notebox}[This is a preview, not the whole story]
We \emph{reduce} distance to LCA and LCA to RMQ here because it is the natural culmination of
the flattening idea. We deliberately do \textbf{not} implement an RMQ or binary lifting in
this chapter -- doing LCA \emph{well} (online, memory-tight, $\bigO{\log n}$ per query via
lifting) is the LCA chapter's job (\chref{c11}); range-minimum \emph{structures} are
\chref{c13}. Take the identity and the tour; the machinery is elsewhere.
\end{notebox}

\subsection{The Algorithm: build the tour and first-occurrence array}

\begin{lstlisting}[caption={The 2n-1 Euler tour, first-occurrence, and distance via LCA}, label={lst:c24-tour-build}]
#include <bits/stdc++.h>
using namespace std;

vector<vector<int>> adj;
vector<int> depth_, first_;
vector<int> tour_;                 // node ids, length 2n-1
vector<int> tourDepth_;            // parallel depths, for the RMQ

void buildTour(int u, int p, int d) {
    first_[u] = (int)tour_.size();  // FIRST time we see u => its index
    depth_[u] = d;
    tour_.push_back(u); tourDepth_.push_back(d);
    for (int v : adj[u]) if (v != p) {
        buildTour(v, u, d + 1);
        tour_.push_back(u); tourDepth_.push_back(d);            // return to u after each child
    }
}
// lca(u,v) = argmin depth over tourDepth in [first[u], first[v]]  -- needs an RMQ (Ch.11/13)
int lcaNaive(int u, int v) {        // O(h) scan, for understanding only
    int l = min(first_[u], first_[v]), r = max(first_[u], first_[v]);
    int best = l;
    for (int i = l; i <= r; ++i) if (tourDepth_[i] < tourDepth_[best]) best = i;
    return tour_[best];
}
long long distUV(int u, int v) {
    int w = lcaNaive(u, v);
    return (long long)depth_[u] + depth_[v] - 2LL * depth_[w];
}
\end{lstlisting}
\begin{complexitybox}
Building $\mathcal{T}$ and first-occurrence is one DFS, $\bigO{n}$ (the tour has $2n-1$
entries). With the $\bigO{h}$ scan above, a query is $\bigO{n}$ worst case -- replace the
scan with a sparse table ($\bigO{n\log n}$ pre, $O(1)$ query) or a segment tree
(\chref{c13}) to make distance queries $O(1)$/$\bigO{\log n}$. Binary lifting
(\chref{c11}) gives $\bigO{\log n}$ LCA \emph{without} the tour. Recursion is $\bigO{h}$
deep -- iterate on a chain.
\end{complexitybox}

\subsection{Worked Example}

\begin{example}{Distance on $G$ via the tour}{c24-dist-run}
$\mathcal{T}=[0,1,4,1,5,1,0,2,0,3,6,3,0]$, so
$\mathrm{first}=[0,1,7,9,2,4,10]$ (by node id) and
$\depth=[0,1,1,1,2,2,2]$. Query $\dist(4,6)$: window
$[\mathrm{first}(4),\mathrm{first}(6)]=[2,10]$, whose minimum-depth entry is node $0$ at
position $6$ (depth $0$) $\Rightarrow \lca(4,6)=0$. Then $\dist(4,6)=\depth4+\depth6-
2\depth0 = 2+2-0 = 4$ (the path $4$--$1$--$0$--$3$--$6$). Query $\dist(1,5)$: window $[1,4]$,
min is node $1$ (depth $1$) $\Rightarrow \lca=1$, $\dist=1+2-2\cdot1=1$.
\end{example}
\begin{tracebox}[Reading $\mathcal{T}$ between first-appearances]
Positions $2..10$ of $\mathcal{T}$ are $[4,1,5,1,0,2,0,3,6]$ with depths
$[2,1,2,1,0,1,0,1,2]$. The smallest depth is $0$, first attained at position $6$ where
$\mathcal{T}[6]=0$; no node deeper than $0$ sits between, and $0$ is the shared ancestor of
$4$ and $6$. Hence $\lca(4,6)=0$, $\dist=4$. \checkmark\ \figref{c24-tour}.
\end{tracebox}

\subsection{Variations and Twists}

\begin{itemize}[nosep]
  \item \textbf{$k$-th node on a path}: walk up from the deeper endpoint to the LCA, then
        down -- $\bigO{\log n}$ with lifting (\chref{c11}).                     \hfill\lv{5}
  \item \textbf{Edge-weighted distance}: replace $\depth$ by root-path \emph{weight}; the
        same $\dist = D_u + D_v - 2D_{\ell}$ holds.                          \hfill\lv{4}
  \item \textbf{Path max/min edge}: stitch two root-paths and combine -- the heavy-light
        add-on (\chref{c16}).                                                  \hfill\lv{5}
  \item \textbf{Diameter via two BFS}: for \emph{unrooted} longest-path, no LCA needed -- but
        leaf-pair counting (LC 1530) uses exactly this combine-at-LCA shape.     \hfill\lv{4}
  \item \textbf{Count pairs within $d$}: DFS returning leaf-distance lists and merging
        (the shape of \probref{c24-leafpairs}).                                 \hfill\lv{4}
\end{itemize}

\subsection{Common Pitfalls}\label{sec:c24-dist-pitfalls}

\begin{warningbox}[Trap 1 -- tour built with the entry-only clock (length $n$, not $2n-1$)]
The depth-minimum LCA property needs the tour that \emph{revisits} a node on every return
from a child. Using the preorder-only array (length $n$) has no re-visits, so the window
between two nodes may not contain their LCA.
\end{warningbox}
\begin{lstlisting}[caption={\textcolor{warningred}{WRONG}: no return steps, too short to hold the LCA}, language=C++, basicstyle=\ttfamily\small]
void walk(int u,int p){ tour_.push_back(u); for(int v:adj[u])if(v!=p)walk(v,u); } // len n
\end{lstlisting}
\begin{lstlisting}[caption={\textcolor{successgreen}{RIGHT}: push the parent again after each child}, language=C++, basicstyle=\ttfamily\small]
void walk(int u,int p){ tour_.push_back(u);
    for(int v:adj[u])if(v!=p){ walk(v,u); tour_.push_back(u);} }   // len 2n-1
\end{lstlisting}

\begin{warningbox}[Trap 2 -- $\min$/$\max$ on first-occurrence before the LCA window]
The window endpoints must be ordered: $[\min(\mathrm{first}\,u,\mathrm{first}\,v),
\max(\dots)]$. Sweeping $[\mathrm{first}\,u, \mathrm{first}\,v]$ blindly is empty (and wrong)
when $u$ is met after $v$.
\end{warningbox}
\begin{lstlisting}[caption={\textcolor{warningred}{WRONG}: window can invert}, language=C++, basicstyle=\ttfamily\small]
int l=first_[u], r=first_[v];          // r may be < l -> loop never runs
\end{lstlisting}
\begin{lstlisting}[caption={\textcolor{successgreen}{RIGHT}: normalise the endpoints}, language=C++, basicstyle=\ttfamily\small]
int l=min(first_[u],first_[v]), r=max(first_[u],first_[v]);
\end{lstlisting}

\begin{warningbox}[Trap 3 -- treating $\mathrm{first}$ as a full LCA implementation]
$O(h)$ window scans are $\bigO{n}$ per query; a batch of $10^5$ distance queries TLEs. The
\emph{scan} is for understanding -- swap in a sparse table / segment tree / binary lifting
(\chref{c11}, \chref{c13}) before using it at scale.
\end{warningbox}
\begin{lstlisting}[caption={\textcolor{warningred}{WRONG}: naive scan under a query load}, language=C++, basicstyle=\ttfamily\small]
for each query: lcaNaive(u,v);        // O(n) each -> O(nq) total, TLE at q,n ~ 1e5
\end{lstlisting}
\begin{lstlisting}[caption={\textcolor{successgreen}{RIGHT}: answer LCA in O(1)/O(log n)}, language=C++, basicstyle=\ttfamily\small]
// RMQ over tourDepth (sparse table / segtree) or binary lifting -- Ch.11/13
\end{lstlisting}

\subsection{Problem Lab}

Two distance problems that avoid needing the RMQ structure (small trees) but lean on the
root-path / LCA \emph{shape}. LC 2096 \emph{Step-By-Step Directions} asks for the moves
between two nodes -- found by comparing their root-paths down to the LCA. LC 1530 \emph{Number
of Good Leaf-Pairs} counts leaf pairs within a distance -- combining leaf depths at their
common ancestor. Both are \thmref{c24-dist} with the LCA located by paths rather than a tour.

% ---------------------------------------------------------------------------
\solvedproblem{c24-pathdir}{Step-By-Step Directions From a Binary Tree Node to Another}{LeetCode 2096}
  {https://leetcode.com/problems/step-by-step-directions-from-a-binary-tree-node-to-another/}{LC Medium L3}
  {root-to-node paths, LCA by trimming common prefix, U/L/R stitching}

\begin{problemstatement}[Step-By-Step Directions]
A \textbf{binary} tree with $n$ nodes valued $1 \dots n$ (each value distinct), given by
\code{root}. Given \code{startValue} and \code{destValue}, find the shortest path from the
node valued \code{startValue} to the node valued \code{destValue} and return it as a string of
moves: \code{'L'} (to left child), \code{'R'} (to right child), \code{'U'} (to parent).
$2 \le n \le 10^5$.
\end{problemstatement}

\begin{attemptgate}[Attempt it yourself first -- \lv{3}, give yourself 20 minutes]
A fair first attempt: find the root-to-node path for each endpoint, then look at where they
diverge.
\peek{the path from $A$ to $B$ goes up to their lowest common ancestor then down -- how do you
find that LCA from the two root-paths?}
\nudge{compute root$\to$start and root$\to$dest as move-strings; strip the shared prefix (the
part to the LCA); the leftover on the start side becomes \code{'U'}s, the leftover on the dest
side is already \code{'L'}/\code{'R'}.}
\fullhint{DFS to record \code{pathStart} and \code{pathDest} (from root, in L/R). Let $i$ be
the first index they differ; LCA is at depth $i$. Answer $=$
\code{string(len(pathStart)}$-i,\text{'U')}$ pathDest.substr($i$). Two DFS passes, no explicit
parent map. Full build below.}
\end{attemptgate}

\solstep{Understanding the Problem}
The unique path from $A$ to $B$ in a tree is: from $A$ up to $\ell=\lca(A,B)$, then down to
$B$. If we know the root-path \emph{move strings} to $A$ and to $B$, their shared prefix is
exactly the root-path to $\ell$; the part of $A$'s path \emph{after} the shared prefix must be
traversed \emph{upward} (one \code{'U'} per edge), and the part of $B$'s path after the prefix
is traversed \emph{downward} (as-is, \code{'L'}/\code{'R'}). This is \thmref{c24-dist}'s
``meet at the LCA'' in move-string form.

\solstep{Key Observations}
\begin{observation}
Root-to-node move strings locate the LCA as the length of the \emph{longest common prefix} --
no parent pointers needed, just two DFS descents that record directions.
\end{observation}
\begin{observation}
$\ell$ at depth $i$ means the start side contributes $(\text{startLen}-i)$ \code{'U'}s and the
dest side contributes its remaining $(\text{destLen}-i)$ moves verbatim.
\end{observation}
\begin{observation}
Recording both paths in one DFS pass (stop when both found) is $\bigO{n}$; the string splice is
$\bigO{h}$.
\end{observation}

\solstep{Think It Through With Me}
\begin{thinkbox}
\deadend{Build a parent map, then walk $A$ and $B$ up until equal.}{Correct, but a
\code{unordered\_map} of parents plus two upward walks is more state than two root-down DFS
paths; and on a skewed $10^5$ tree the upward walks are fine, but the map is pure overhead.}
\deadend{Do a full BFS from $A$ to $B$ over the tree with \code{U/L/R} edges.}{You must model
the parent as an explicit edge (turning the binary tree into an undirected graph) -- heavier
than exploiting the tree's unique-path / LCA structure directly.}
\aha{DFS for the two root-path move strings; trim their common prefix; \code{'U'} the rest of
the start side, append the rest of the dest side. That \emph{is} the $A$-up-to-LCA-down-to-$B$
path.}
\end{thinkbox}

\solstep{Solution}
Two DFS (or one) that record the direction string from the root to the node with value
\code{startValue} and to \code{destValue}. Advance $i$ while both strings agree. Return
\code{string(start.size()-i, 'U') + dest.substr(i)}.

\solstep{Proof of Correctness}
\begin{proofbox}
The root-to-node direction string is unique per node (a tree has one root-path, each step
left/right by \chref{c01} uniqueness). Its length is $\depth$. Two such strings agree for
exactly the first $\depth(\ell)$ characters, where $\ell=\lca(A,B)$, because the root-paths
coincide precisely down to their lowest shared node and then diverge; so the first differing
index is $i=\depth(\ell)$. From $A$ the path to $\ell$ reverses $A$'s last
$\depth(A)-\depth(\ell)=|start|-i$ moves, each a \code{'U'}. From $\ell$ to $B$ follows $B$'s
last $|dest|-i$ moves unchanged. Concatenation is the unique simple $A$--$B$ path, hence the
shortest. \thmref{c24-dist} is consistent: the returned length
$(|start|-i)+(|dest|-i)$ equals $\depth A + \depth B - 2\depth\ell$.
\end{proofbox}

\solstep{Implementation}
\begin{lstlisting}[caption={Step-By-Step Directions via root-path LCA (LeetCode 2096)}, label={lst:c24-pathdir}]
#include <bits/stdc++.h>
using namespace std;

struct TreeNode { int val; TreeNode *left, *right;
    TreeNode(int v):val(v),left(nullptr),right(nullptr){} };

class Solution {
public:
    // find direction string from u down to target, accumulating into path
    bool dfs(TreeNode* u, int target, string& path) {
        if (!u) return false;
        if (u->val == target) return true;
        path.push_back('L');
        if (dfs(u->left, target, path)) return true;
        path.pop_back();
        path.push_back('R');
        if (dfs(u->right, target, path)) return true;
        path.pop_back();
        return false;
    }
    string getDirections(TreeNode* root, int startValue, int destValue) {
        string s, d;
        dfs(root, startValue, s);
        dfs(root, destValue, d);
        int i = 0;
        while (i < (int)s.size() && i < (int)d.size() && s[i] == d[i]) ++i;  // common prefix
        return string(s.size() - i, 'U') + d.substr(i);
    }
};
\end{lstlisting}

\solstep{Trace}
\begin{tracebox}[$R$, \figref{c03-running}: start$=5$, dest$=3$]
Root$\to$5 path $=$ \code{"LR"} ($1$'s left $2$, then right $5$). Root$\to$3 path $=$
\code{"R"}. Common prefix length $i=0$ (\code{L}$\ne\ $\code{R}), so $\ell=$ root $=1$.
Result $=$ \code{"UU"} (two ups from $5$ to $1$) $+$ \code{"R"} (down to $3$) $=$
\code{"UUR"}. Length $2+1=3=\depth5+\depth3-2\depth1=2+1-0$ \checkmark\
\probref{c24-pathdir}.
\end{tracebox}

\solstep{Complexity}
\begin{complexitybox}
Two DFS descents worst $\bigO{n}$ each; splice $\bigO{h}$. Time $\bigO{n}$, space
$\bigO{h}$ for the recursion plus the path strings. On a $10^5$ skewed tree the DFS is
$\bigO{n}$ deep -- make it iterative. Verdict: \ac.
\end{complexitybox}

\solstep{Pitfalls and Variants}
\begin{warningbox}
When $\ell$ \emph{is} $A$ or $B$ the common prefix can consume one entire string; the loop
guards both bounds so $i$ stops correctly. Also do not forget to \code{pop\_back} on the
failed \code{'R'} branch -- an un-popped tail corrupts the dest string. Variant: with parent
pointers this becomes the two-upward-walk method; equivalent, more state.
\end{warningbox}

\begin{insightbox}[Key takeaways]
\begin{enumerate}[nosep, leftmargin=1.6em]
  \item The LCA $=$ longest common prefix of two root-path direction strings -- locating it
        \emph{is} the ``meet at LCA'' of \thmref{c24-dist}.
  \item Any ``path between two nodes'' answer is up-to-LCA (as \code{U}) then down-from-LCA
        (as recorded).
  \item Backtracking must undo \emph{every} push; a leaked branch corrupts the result.
\end{enumerate}
\end{insightbox}

% ---------------------------------------------------------------------------
\solvedproblem{c24-leafpairs}{Number of Good Leaf Nodes Pairs}{LeetCode 1530}
  {https://leetcode.com/problems/number-of-good-leaf-nodes-pairs/}{LC Medium L4}
  {combine leaf-depth lists at LCA, count pairs within d, small n}

\begin{problemstatement}[Number of Good Leaf Nodes Pairs]
Given a \textbf{binary} tree \code{root} and an integer \code{distance}. A leaf is a node with
no children. Count \textbf{pairs of different leaf nodes} whose shortest path between them has
length $\le$ \code{distance}. $2 \le n \le 2000$, $1 \le \text{distance} \le 10$.
\end{problemstatement}

\begin{attemptgate}[Attempt it yourself first -- \lv{4}, give yourself 30 minutes]
A fair first attempt: think about where the path between two leaves ``bends'' -- their LCA --
and what you would want each subtree to report.
\peek{at a node $u$, a leaf-pair whose LCA is $u$ has one leaf in $u$'s left and one in its
right; what distances should the children hand back?}
\nudge{each DFS returns the multiset of \emph{distances from $u$ down to the leaves of
$\subtree(u)$}; at $u$, pair every left distance $a$ with every right distance $b$ and count
those with $a+b \le$ \code{distance}; then return the child distances, each $+1$.}
\fullhint{Bottom-up: \code{dfs(u)} returns a list of distances from $u$ to leaves in its
subtree (leaf $\Rightarrow$ empty list handled as the node itself). For internal $u$:
$\mathit{lst}=\text{left}{+}1$ merged with $\text{right}{+}1$; before merging, cross-count
pairs $(a\in\text{left}, b\in\text{right})$ with $a+b\le d$. $n\le2000$, distance$\le10$, so
$O(n\cdot \text{leaves})$ or a small array of size $d$ is fine. Full build below.}
\end{attemptgate}

\solstep{Understanding the Problem}
Every leaf pair has a unique LCA $\ell$; the path length is $\dist(\text{leaf}_1,
\text{leaf}_2) = a + b$ where $a, b$ are the two leaves' distances to $\ell$. So if we
process \emph{at} each node the leaves reachable within $\le d$ below it, partitioned by which
child they descend from, and count left$\times$right pairs summing to $\le d$, we count every
good pair \emph{exactly once} (at its LCA). A bottom-up DFS returns, per node, the distances to
its descendant leaves.

\solstep{Key Observations}
\begin{observation}
A leaf pair is counted at its LCA, splitting into one leaf strictly in the left subtree and one
in the right (if both are on the same side, their LCA is lower). This gives exact-once counting.
\end{observation}
\begin{observation}
Because $\text{distance} \le 10$, we only ever need to propagate distances $\le d$; anything
deeper can be pruned -- a per-node array \code{cnt[1..d]} (leaves at each distance) keeps
merges cheap.
\end{observation}
\begin{observation}
With $n \le 2000$, even the plain ``return a vector of leaf distances, cross-multiply $O(
\mathit{left}\cdot\mathit{right})$'' version passes; the $\bigO{n\cdot d}$ histogram is the
polished form.
\end{observation}

\solstep{Think It Through With Me}
\begin{thinkbox}
\deadend{Enumerate all leaf pairs and run a BFS/LCA for each distance.}{$\binom{L}{2}$ leaf
pairs, each needing an LCA -- with $L$ up to $\sim1000$ that's $\sim5\cdot10^5$ queries; fine
at $n\le2000$ but it re-derives distances without structure. The bottom-up combine counts the
same thing with one DFS.}
\deadend{Count pairs whose path passes through $u$ using \emph{all} leaves in
$\subtree(u)$.}{That double-counts: a pair whose LCA is \emph{below} $u$ also lies in $u$'s
subtree. Only pairs split across $u$'s \emph{two different children} have LCA exactly $u$.}
\aha{At each node cross-count left-leaves $\times$ right-leaves with $a+b\le d$ (LCA $=u$),
return the merged $+1$ distance list -- each pair counted once, at its true LCA.}
\end{thinkbox}

\solstep{The Approach Ladder}
\begin{center}\renewcommand{\arraystretch}{1.3}
\begin{tabular}{@{}c p{0.235\linewidth} p{0.28\linewidth} l c@{}}
\toprule
\textbf{\#} & \textbf{Approach} & \textbf{Idea} & \textbf{Time} & \textbf{Verdict} \\
\midrule
1 & \lv{3} enumerate leaf pairs + LCA & dist formula per pair & $O(L^2\log n)$ & \ac \\
2 & \lv{4} bottom-up distance lists & cross-count at each LCA & $O(n\cdot L)$ & \ac \\
3 & \lv{4} histogram capped at $d$ & \code{cnt[1..d]} merge & $O(n\cdot d)$ & \ac \\
\bottomrule
\end{tabular}\end{center}
Rung 3 exploits distance $\le 10$; rungs 1--2 are conceptually simpler and pass at $n\le2000$.

\solstep{Solution}
Bottom-up DFS returning \code{vector<int>} of distances to descendant leaves $\le d$. For a
leaf, return $\{0\}$ at itself (handled by parent). For an internal node: get
\code{left}, \code{right} (already $+1$-ed by the caller or here), cross-count all
$(a,b)$ with $a+b \le d$ into the answer, then return the concatenation. Incremental $+1$
as we ascend.

\solstep{Proof of Correctness}
\begin{proofbox}
\emph{Claim:} the DFS returns, at node $u$, the multiset of distances $\dist(u, \text{leaf})$
for leaves in $\subtree(u)$ (those $> d$ discarded safely since they can form no good pair with
any non-negative partner). A leaf trivially has distance $0$ to itself; ascending adds $1$ per
edge. \emph{Counting:} a good pair $(x,y)$ of distinct leaves has a unique LCA $\ell$; $x$
lies strictly in one child subtree of $\ell$ and $y$ in the other (else $\ell$ was not lowest).
At $\ell$ the algorithm forms all cross pairs $(a\in\text{left}, b\in\text{right})$ and counts
exactly those with $a+b \le d$; but $a+b = \dist(\ell,x)+\dist(\ell,y) = \dist(x,y)$
(\thmref{c24-dist} with $\depth\ell$ cancelling). Thus $(x,y)$ is counted at $\ell$ iff good,
and at no other node (the same-side pairs are excluded). Hence each good pair contributes
exactly $1$ and each counted pair is good.
\end{proofbox}

\solstep{Implementation}
\begin{lstlisting}[caption={Good Leaf Pairs by cross-counting at the LCA (LeetCode 1530)}, label={lst:c24-leafpairs}]
#include <bits/stdc++.h>
using namespace std;

struct TreeNode { int val; TreeNode *left, *right;
    TreeNode(int v):val(v),left(nullptr),right(nullptr){} };

class Solution {
public:
    int ans = 0;
    int D;
    // returns distances from u to leaves in subtree(u), capped at D
    vector<int> dfs(TreeNode* u) {
        if (!u) return {};
        if (!u->left && !u->right) return {0};          // leaf: distance 0 to itself
        vector<int> L = dfs(u->left),  R = dfs(u->right);
        for (int& a : L) ++a;                            // child's leaves are one deeper
        for (int& b : R) ++b;
        for (int a : L) for (int b : R) if (a + b <= D) ++ans;   // cross pairs: LCA = u
        vector<int> out;
        for (int a : L) if (a <= D) out.push_back(a);    // propagate upward
        for (int b : R) if (b <= D) out.push_back(b);
        return out;
    }
    int countPairs(TreeNode* root, int distance) {
        D = distance; ans = 0; dfs(root); return ans;
    }
};
\end{lstlisting}
\noindent The two \code{++} loops bind by \textbf{reference} (\code{int\&}); a by-value
range-\code{for} would be a silent no-op (Pitfalls, trap (i)).

\solstep{Trace}
\begin{tracebox}[$R$, \figref{c03-running}: leaves $4,5,6$, distance $=3$]
Node $2$: children $4,5$ (leaves, return $\{0\}$ each); $+1$ $\Rightarrow$ $L=\{1\},R=\{1\}$;
cross pair $1+1=2\le3$ $\Rightarrow$ count $1$ (pair $4,5$); return $\{1,1\}$.
Node $3$: child $6$ leaf; after $+1$, one side $\{1\}$, other empty; count $0$; return $\{1\}$.
Node $1$: $L$ (from $2$) $=\{1,1\}+1=\{2,2\}$, $R$ (from $3$) $=\{1\}+1=\{2\}$; cross pairs:
$2+2=4>3$ twice $\Rightarrow$ add $0$; total stays $\mathbf{1}$. \checkmark\
\probref{c24-leafpairs}.
\end{tracebox}

\solstep{Complexity}
\begin{complexitybox}
Each node merges its two children's (capped) distance lists: the naive cross product is
$\bigO{L^2}$ at a node worst case, total $\bigO{n\cdot L}$ with the $d$-cap since lists
hold $\le d$ distinct entries ($\bigO{n\cdot d^2}$ worst). At $n\le2000$, $d\le10$ this is
trivially fast; the histogram form is $\bigO{n\cdot d}$. Space $\bigO{h\cdot d}$. Verdict: \ac.
\end{complexitybox}

\solstep{Pitfalls and Variants}
\begin{warningbox}
Two real bugs: (i) the copy-vs-reference increment above (a range-\code{for (int a : L)} does
not modify $L$); (ii) counting \emph{all} pairs in $\subtree(u)$ rather than only
left$\times$right cross pairs, which double-counts pairs whose LCA is below $u$. Variant:
swap the list for a size-$d$ histogram \code{cnt[]} and merge in $\bigO{d}$ with a convolution
for the cross count.
\end{warningbox}

\begin{insightbox}[Key takeaways -- three things to remember]
\begin{enumerate}[nosep, leftmargin=1.6em]
  \item ``Count leaf pairs within $d$'' $=$ combine at the \emph{LCA}, counting only
        cross-child pairs -- the exact-once trick \thmref{c24-dist} underwrites.
  \item Returning a (capped) distance multiset bottom-up is the distance-indexed DP pattern
        the DP chapters generalise.
  \item Range-\code{for} over a value (not \code{\&}) is a silent no-op -- a whole class of
        WA.
\end{enumerate}
\end{insightbox}

\begin{hardeningbox}
\begin{itemize}[nosep, leftmargin=1.4em]
  \item \textbf{Big $n$, small $d$}: use the $\bigO{d}$ histogram merge; the list-of-leaves
        version is $\bigO{L}$ per node and can $\bigO{n^2}$ if $d$ were large.
  \item \textbf{Depth}: recursion $\bigO{h}$ -- at $n=2000$ fine; generalise with an explicit
        stack if $n$ grows.
\end{itemize}
\end{hardeningbox}

\subsection{Practice Problems}

\begin{stuckbox}
\begin{itemize}[nosep, leftmargin=1.4em]
  \item Does your tour push the parent back \emph{after each child} (length $2n-1$)? (Trap 1.)
  \item Are the RMQ window endpoints $\min/\max$ of first-occurrences? (Trap 2.)
  \item Is the LCA-locating \emph{scan} a stand-in you would replace by a real structure at
        query load? (Trap 3, \chref{c11}.)
  \item In a pair count, are you only pairing \emph{across} children (LCA $=u$), not within one
        side? (\probref{c24-leafpairs}.)
\end{itemize}
\end{stuckbox}

\begin{practiceset}
  \practice{Diameter of N-Ary Tree}{LeetCode 1522}
    {https://leetcode.com/problems/diameter-of-n-ary-tree/}{LC Medium L3}
    {The \emph{length} of the longest path between any two nodes (in edges) -- the maximum of
    $\dist(u,v)$ over all pairs. The path's top is some node $a$ (its LCA), and its length is
    the two tallest downward chains hanging below $a$. Return, from each node bottom-up, the
    tallest child height; the best path through $a$ is the sum of its \emph{two} tallest child
    chains (\thmref{c24-dist} with the common ancestor fixed at $a$). Track the global max while
    returning $1+\max$ child height upward.}
\end{practiceset}

\begin{recallbox}[10-second refresher -- 5.4 Distance \& RMQ Preview]
The $2n{-}1$ tour $\mathcal{T}$ records the node on every arrival \emph{and} return; index it by
$\mathrm{first}[v]=\tin(v)$. Then $\dist(u,v)=\depth u+\depth v-2\depth{\lca}$
(\thmref{c24-dist}), and $\lca$ is the \emph{min-depth} node in $\mathcal{T}$ between the two
firsts -- an RMQ (\propref{c24-rmqpreview}), \emph{proved and built} in \chref{c11}. Paths to
the LCA also answer directions (\probref{c24-pathdir}) and leaf-pairs-by-distance
(\probref{c24-leafpairs}).
\selftest{
  \item Write the DFS that emits the $2n{-}1$ tour and the first-occurrence array
        (\coderef{c24-tour-build}).
  \item State and prove $\dist(u,v)=\depth u+\depth v-2\depth{\lca}$.
  \item Explain why the min-depth node between two firsts is the LCA (and where the full proof
        lives).
  \item On $G$, compute $\dist(4,6)$ from the tour window and $\depth$.
  \item For \probref{c24-leafpairs}, say why a pair is counted at its LCA and only across
        children.
}
\end{recallbox}
